\chapter{Angina}
\index{Angina}

\section*{Definition and Overview}
Angina is a clinical syndrome of chest pain or discomfort caused by a temporary reduction in the supply of oxygen-rich blood to the heart muscle (myocardial ischaemia), most commonly due to narrowing of the coronary arteries by atherosclerotic plaques. When the heart works harder -- during exertion, emotional stress, cold weather, or after a heavy meal -- the narrowed arteries cannot deliver enough blood, and the heart muscle responds with a characteristic squeezing, heavy, pressing, or tight discomfort.

Angina is not a heart attack, but it is an important warning sign of coronary heart disease. Stable angina is predictable, occurs on exertion, and settles quickly with rest or glyceryl trinitrate (GTN). Unstable angina is a more serious condition in which symptoms are new, occur at rest, or increase in severity and frequency; it is an acute coronary syndrome and can herald a heart attack, requiring urgent medical assessment.

\section*{Epidemiology}
Coronary heart disease is one of the leading causes of death and disability worldwide, and angina is one of its most common presentations. The prevalence of angina increases with age, and it is more common in men at younger ages, with the sex gap narrowing after menopause. Major modifiable risk factors include smoking, hypertension, elevated cholesterol, diabetes, physical inactivity, obesity, and excessive alcohol consumption; non-modifiable factors include age, male sex, and a family history of premature coronary disease.

Rates of angina and of deaths from coronary heart disease have fallen in many developed countries over recent decades because of better prevention and treatment, but the absolute burden remains high, particularly in low- and middle-income countries and in disadvantaged groups. A substantial proportion of people with angina are unaware that their symptoms are cardiac in origin, which delays diagnosis and treatment.

\section*{Aetiology and Pathophysiology}
Angina results from an imbalance between myocardial oxygen demand and coronary blood supply. The main mechanisms include:

\begin{itemize}
    \item \textbf{Atherosclerosis:} fatty and fibrous plaques accumulate in the coronary artery walls, narrowing the lumen and limiting blood flow; this underlies most cases
    \item \textbf{Coronary artery spasm:} a temporary intense constriction of an artery reduces flow even in the absence of significant plaque (variant or Prinzmetal angina)
    \item \textbf{Increased myocardial oxygen demand:} exercise, emotion, cold weather, and heavy meals increase the heart's workload beyond what the narrowed vessels can supply
    \item \textbf{Coronary microvascular dysfunction:} disease of the small resistance vessels can cause angina (often in women) even when the large coronary arteries appear normal at angiography
    \item \textbf{Risk factors:} smoking, hypertension, diabetes, hyperlipidaemia, and a family history damage the arterial endothelium, promote inflammation, and accelerate plaque formation
\end{itemize}

When demand exceeds supply, the heart muscle becomes transiently ischaemic. The ischaemia activates nerve endings, producing referred discomfort, and can impair contraction of the affected muscle and disturb the heart's electrical rhythm.

\section*{Clinical Features}
\begin{itemize}
    \item \textbf{Character:} a heavy, squeezing, pressing, or tight sensation in the centre of the chest, often described as a band or weight
    \item \textbf{Radiation:} discomfort spreading to the left arm, shoulder, jaw, neck, or back
    \item \textbf{Precipitants:} exertion, emotional stress, cold weather, or a large meal
    \item \textbf{Relief:} symptoms usually settle within a few minutes of rest or GTN
    \item \textbf{Associated symptoms:} breathlessness, fatigue, nausea, sweating, or a feeling of indigestion, which may be the dominant feature, especially in women, older people, and people with diabetes
    \item \textbf{Stable angina:} a predictable pattern of symptoms occurring over weeks or months
    \item \textbf{Unstable angina:} new, more frequent, more severe, or rest-related symptoms, or symptoms that no longer respond to GTN
\end{itemize}

Some people, particularly older adults and those with diabetes, experience "silent" ischaemia or atypical symptoms such as isolated breathlessness, and this can delay recognition.

\section*{Differential Diagnosis}
Chest pain has many causes, and the following should be distinguished from angina:

\begin{itemize}
    \item Acute myocardial infarction and other acute coronary syndromes
    \item Gastro-oesophageal reflux disease (heartburn) and oesophageal spasm
    \item Musculoskeletal pain: costochondritis, intercostal muscle strain, or rib injury
    \item Pericarditis and other structural heart conditions
    \item Pulmonary causes: pneumonia, pleurisy, pulmonary embolism, and asthma
    \item Anxiety, panic attacks, and hyperventilation
    \item Vascular emergencies such as aortic dissection
\end{itemize}

The history is the most important diagnostic tool: classic exertional central chest discomfort relieved by rest is highly suggestive of angina, whereas sharp, localised, position-dependent, or brief stabbing pain usually points away from it.

\section*{When to Seek Medical Care}
People should see a doctor without delay if they experience chest discomfort with exertion that they have not had before, if known angina becomes more frequent or severe, or if symptoms no longer settle with rest and their usual medicines. Anyone with risk factors for heart disease and possible angina should be assessed promptly, since early treatment reduces the risk of a heart attack.

\textbf{Contact your local emergency services for:}
\begin{itemize}
    \item Chest pain that lasts more than a few minutes or does not settle with rest or GTN
    \item Chest pain occurring at rest, or that is new, sudden, or more severe than usual
    \item Chest pain with breathlessness, sweating, nausea, vomiting, or collapse
    \item Pain spreading to the arm, jaw, neck, or back, or a feeling of faintness
\end{itemize}

\section*{Diagnosis and Investigations}
Diagnosis combines clinical assessment with investigations to confirm ischaemia and assess risk:

\begin{itemize}
    \item \textbf{History and examination:} detailed description of the pain -- its site, character, radiation, precipitants, and relieving factors -- plus a cardiovascular examination and measurement of blood pressure and pulse
    \item \textbf{Resting ECG:} may be normal between episodes; an ECG recorded during pain is valuable if it shows ischaemic changes
    \item \textbf{Blood tests:} high-sensitivity troponin during symptoms to help exclude a heart attack, together with cholesterol, blood glucose or HbA1c, and kidney function
    \item \textbf{Functional testing:} exercise stress ECG, or stress imaging such as stress echocardiography or myocardial perfusion scanning, to demonstrate inducible ischaemia
    \item \textbf{CT coronary angiography:} a non-invasive scan that can exclude significant coronary narrowing
    \item \textbf{Invasive coronary angiography:} the definitive anatomical test, used when non-invasive testing is inconclusive or high risk is suspected
    \item \textbf{Risk scoring:} validated tools (such as the GRACE score in acute settings) to guide further investigation and treatment
\end{itemize}

\section*{Management}
Management aims to relieve symptoms, slow disease progression, and prevent heart attacks:

\begin{itemize}
    \item \textbf{Lifestyle modification:} smoking cessation, a heart-healthy diet, regular physical activity, weight loss if needed, and stress management
    \item \textbf{GTN:} a sublingual spray or tablet to relieve an acute attack and for use before exertion
    \item \textbf{Antiplatelet therapy:} low-dose aspirin (and sometimes a second agent after an acute event) to reduce clot formation
    \item \textbf{Statins:} to lower cholesterol and stabilise plaques
    \item \textbf{Anti-anginal drugs:} beta-blockers or calcium channel blockers as first-line treatment to reduce myocardial oxygen demand; long-acting nitrates, ivabradine, or ranolazine are options in some patients
    \item \textbf{Risk factor control:} aggressive treatment of hypertension and diabetes
    \item \textbf{Revascularisation:} percutaneous coronary intervention (PCI) with stenting, or coronary artery bypass grafting (CABG) for extensive disease, offered where symptoms are uncontrolled or risk is high
    \item \textbf{Cardiac rehabilitation:} a supervised programme of exercise, education, and psychological support
\end{itemize}

\section*{Complications}
\begin{itemize}
    \item Myocardial infarction if a vulnerable plaque ruptures and occludes an artery
    \item Progression from stable to unstable angina and acute coronary syndrome
    \item Arrhythmias, including ventricular tachycardias
    \item Heart failure from recurrent ischaemic damage to the heart muscle
    \item Sudden cardiac death if ischaemia triggers a life-threatening rhythm
    \item Complications of treatment, including bleeding from antiplatelet agents and the small procedural risks of angiography and revascularisation
\end{itemize}

\section*{Prognosis and Outcomes}
With modern treatment, most people with stable angina live for many years with good symptom control and a low annual risk of heart attack. Angina is usually a long-term condition, but lifestyle change and medication can slow the progression of atherosclerosis and improve the outlook substantially. Unstable angina carries a higher short-term risk and requires urgent assessment and treatment. Smoking cessation and control of blood pressure and cholesterol are among the most powerful measures for improving prognosis, and outcomes are best when people adhere to treatment and attend cardiac rehabilitation.

\section*{Prevention and Self-Care}
\begin{itemize}
    \item Do not smoke, and avoid exposure to second-hand smoke
    \item Keep blood pressure, cholesterol, and blood sugar within healthy ranges
    \item Eat a diet rich in vegetables, fruit, whole grains, and fish, and limit salt, saturated fat, and added sugar
    \item Be physically active on most days, within a plan agreed with your doctor
    \item Maintain a healthy weight and limit alcohol
    \item Take prescribed medicines exactly as directed and always carry your GTN spray or tablets
    \item Learn the warning signs of a heart attack and act quickly if they occur
    \item Attend cardiac rehabilitation and scheduled follow-up appointments
\end{itemize}

\section*{Sources and Further Reading}
The following sources provide reliable further information on angina:

\begin{itemize}
    \item NHS. \textit{Overview of angina, its symptoms, causes, and treatment.} https://www.nhs.uk/conditions/angina/
    \item Mayo Clinic. \textit{Information on the types, diagnosis, and management of angina.} https://www.mayoclinic.org/diseases-conditions/angina/symptoms-causes/syc-20369373
    \item British Heart Foundation. \textit{Plain-language information on angina and living with it.} https://www.bhf.org.uk/informationsupport/heart-matters
    \item NICE. \textit{Clinical guideline on the assessment and diagnosis of chest pain of recent onset.} https://www.nice.org.uk/guidance/cg95
    \item MSD Manual. \textit{Angina pectoris in the professional health section.} https://www.msdmanuals.com/professional/cardiovascular-disorders/coronary-artery-disease/angina-pectoris
\end{itemize}
